\begin{enumerate}
\item[7.1] \thm{$\hat{T} = e^{-b\hat{D}}$}
  \begin{proof}
    \begin{align*}
      \bra{p} \exp(-b\hat{D}) &= \bra{p} \exp(-bD(p))
    \end{align*}
    where $D(p) = ip/\hbar$ is the eigenvalue of $\hat{D}$ for $\bra{p}$, and knowing that
    \begin{equation}
      \bra{p} \hat{T} = \exp(\frac{-ibp}{\hbar}) \bra{p}
    \end{equation}
    from problem 3.1,
    \begin{align*}
      \bra{p} \hat{T} = \exp(-bD(p)) \bra{p} = \exp(-b\hat{D}) \bra{p},
    \end{align*}
    thus
    \begin{align*}
      \hat{T} = e^{-b\hat{D}}.
    \end{align*}
  \end{proof}
\item[7.2] \thm{$\bra{x} \hat{p} \ket{\psi} = \bra{x} \hat{p} \hat{T} \ket{\psi}$}
  \begin{proof}
    Given that $\hat{T} = e^{-b\hat{D}}$ per problem 7.1, $\braket{x}{\psi}$ and $\bra{x} \hat{T} \ket{\psi}$ differ by only the phase $e^{-b\hat{D}}$ which cannot have any effect on the probability as $\forall \theta, |e^{i\theta}|^2 = 1$.
  \end{proof}
\item[7.3] \thm{$\forall \psi, \forall t$, $|\braket{E}{\psi}|^2 = |\bra{E} \hat{U}_t \ket{\psi}|^2$.}
  \begin{proof}
    $\hat{U}_t = \exp(\frac{-it\hat{H}}{\hbar})$ is a unit vector on the complex plane, so it cannot have any effect on the probability.
  \end{proof}
\item[7.4] What is the characteristic frequency $\Omega/2\pi$ for a neutron in a magnetic field of $\SI{1}{\kilo\gauss}$? In what region of the electromagnetic spectrum is this: radio, microwave, infrared, visible, or x-ray?

  The characteristic frequency is
  \begin{equation*}
    \omega = \frac{B}{2\pi} \abs{\Omega} = \frac{B}{2\pi} \abs{\frac{g_n e}{2 m_n}}
    = \frac{\SI{1E3}{\kilo\gauss}}{2\pi}
    \abs{\frac{(-3.83)(\SI{1.60E-19}{\coulomb})}
         {2(\SI{1.67E-27}{\kilo\gram})}}
    \simeq \SI{2.92}{\mega\hertz}
  \end{equation*}
  and this is a radiowave.
\item[7.5] \thm{$\hat{U}_t \ket{\mu = 1/2} = \frac{\hat{U}_t}{\sqrt{2}}(\ket{m = 1/2} + i\ket{m = -1/2})$ where $\mu$ and $m$ are eigenvalues of $\hat{S}_y$ and $\hat{S}_z$ respectively.}
  \begin{proof}
    The unit time transform can be removed from both sides, so the proof is just
    $$\ket{\mu = 1/2} = \ket{m = 1/2} + i\ket{m = -1/2}$$
    and this can be proven using the change-of-basis matrix from $\eig \hat{S}_y$ to $\eig \hat{S}_z$. This can be produed using the spin-$1/2$ rotation operator 
  \end{proof}
\end{enumerate}
